\INDEX{證明：}\setcounter{equation}{0}
因為 $S$ 在射線 $GC$ 上，所以存在正數 $\lambda > 0$ 使得 $\overrightarrow{GS}=\lambda\overrightarrow{GC}$。
記 $AB$ 的中點為 $M$，由重心定理可知 $\overrightarrow{GM}=-\dfrac{1}{2}\overrightarrow{GC}$，所以有
\begin{equation}
  \overrightarrow{GG}=\overrightarrow{0}=\frac{1}{1+2\lambda}\overrightarrow{GS}+\frac{2\lambda}{1+2\lambda}\overrightarrow{GM}
\end{equation}
由定比分點的性質可知，$G$ 在線段 $MS$ 上，進一步可得 $G$ 在 $\angle{MBS}=\angle{ABS}$ 的內部。

由題中的已知條件可知
\begin{equation}
  \angle{ABS}=180^{\circ}-\angle{BGC}=180^{\circ}-\angle{BGS}
\end{equation}
結合「$G$ 在 $\angle{ABS}$ 的內部」的推論，利用弦切角定理的逆定理可得出 $AB$ 是 $\triangle{BGS}$ 外接圓的切線。

結合 $M$ 點的中點條件，由切割線定理可知
\begin{equation}
  \overrightarrow{MA}^2=\overrightarrow{MB}^2=\overrightarrow{MG}\cdot\overrightarrow{MS}
\end{equation}
所以又可得知 $MA$ 是 $\triangle{AGS}$ 外接圓的切線，即 $AB$ 是 $\triangle{AGS}$ 外接圓的切線，進一步可得
\begin{equation}
  \angle{BAG}=\angle{MAG}=\angle{ASG}\label{eq:p5_s1}
\end{equation}
同理可得 $AC$ 是 $\triangle{AGR}$ 外接圓及 $\triangle{CGR}$ 外接圓的外公切線，以及 $\angle{CAG}=\angle{ARG}$。

最終可知
\begin{align}
  \angle{BGC}&=\angle{RGS}\notag\\
  &=\angle{RAS}+\angle{ARG}+\angle{ASG}\notag\\
  &=\angle{RAS}+\angle{CAG}+\angle{BAG}\notag\\
  &=\angle{RAS}+\angle{BAC}
\end{align}